\chapter{Installation}
\label{ch:installation}

\newthought{This chapter is for anyone setting up \Jtwo} --- on a laptop, a group
server, or a cluster login node. It requires nothing beyond a shell and Python. After
following it, you have a working \cli{Jarvis} command, a running Redis service, and a
verified installation --- ready for the quick start in Chapter~\ref{ch:quickstart}.

The whole installation is two components:

\begin{enumerate}
  \item the \textbf{Python package} \code{jarvishep2} --- the runtime, the
        \cli{Jarvis} command, and its three core plugin
        packages;\sidenote{The import name and the \textsc{pypi} distribution name
        differ: \code{import jarvishep2} in Python, but \cli{pip install Jarvis-HEP}
        and \cli{pip show Jarvis-HEP} at the shell. Section~\ref{sec:uninstall} and
        Chapter~\ref{ch:reproducibility} use the distribution name for exactly this
        reason.}
  \item a local \textbf{Redis service} --- the small data server every scan
        coordinates through (Section~\ref{sec:runtime-picture}).
\end{enumerate}

\section{Check the requirements}
\label{sec:requirements}

\begin{itemize}
  \item \textbf{Python 3.10 or newer}, on macOS or Linux. Check with
        \cli{python3 --version}.
  \item \textbf{Redis} for actual scans; Section~\ref{sec:redis-install} installs it.
        The test suite does \emph{not} need it --- tests run on a stand-in
        (\code{fakeredis}).
\end{itemize}

\begin{jnote}
\Jtwo\ always connects to the local Redis at \code{127.0.0.1:6379}. This is fixed by
design --- there is no task-card or CLI setting for it, so there is nothing to
configure and nothing to get wrong.
\end{jnote}

\section{Install the package}
\label{sec:pip-install}

No source checkout, no cloning --- one command, straight from \textsc{pypi}:

\begin{enumerate}
  \item Install:
\begin{terminal}
\begin{jcmd}
python3 -m pip install Jarvis-HEP
\end{jcmd}
\end{terminal}
  \item Confirm the command is on your \code{PATH}; it prints the usage text:
\begin{terminal}
\begin{jcmd}
Jarvis --help
\end{jcmd}
\end{terminal}
\end{enumerate}

Three plugin packages install automatically as core dependencies, with no extra to
ask for: \Portal\ (calculator file formats, Section~\ref{sec:format-catalog}),
\Operas\ (the operator registry, Chapter~\ref{ch:jarvis-operas}), and \JPlot\
(Chapter~\ref{ch:jarvisplot}). This already gives you \cli{Jarvis plot},
\cli{Jarvis portal}, and \cli{Jarvis operas}.

\begin{jwarn}
What it does \emph{not} give you yet is \code{redis}: a real, Redis-backed scan needs
three more packages installed alongside it (Table~\ref{tab:extras}). \cli{Jarvis man},
\cli{Jarvis plot}, and \cli{Jarvis validate} all work without them; \cli{Jarvis run}
and \cli{Jarvis check} do not.
\end{jwarn}

\begin{table}[ht]
  \footnotesize
  \begingroup
  \setlength{\tabcolsep}{0pt}
  \begin{tabular}{@{}p{0.34\linewidth}p{0.66\linewidth}@{}}
    \toprule
    \textbf{For} & \textbf{Also install} \\
    \midrule
    real, Redis-backed scans & \code{redis}, \code{msgpack}, \code{aiofiles} \\
    running the test suite   & \code{pytest}, \code{fakeredis}, \code{colorlog} \\
    \bottomrule
  \end{tabular}
  \endgroup
  \caption{Additional packages by purpose. Name as many as you need on the same
  line: \cli{pip install Jarvis-HEP redis msgpack aiofiles}.}
  \label{tab:extras}
\end{table}

\begin{jnote}
\Jtwo\ is its own distribution (\code{Jarvis-HEP}) with its own command
(\cli{Jarvis}). It coexists in one environment with a version~1 installation ---
neither package touches the other.
\end{jnote}

\section{Check the version}
\label{sec:verify-version}

\begin{terminal}
\begin{jcmd}
Jarvis --version    # or the short form: Jarvis -v
\end{jcmd}
\end{terminal}

\noindent This prints the \Jarvis\ banner with the runtime version of your installed
package on its last line (Figure~\ref{fig:version-banner}) --- the fastest way to
confirm exactly which release you are running, without touching Redis or a task card.
The coloured dot-matrix on the left is the \Jarvis\ icon; the wordmark fades from blue
to gold, and the version line is read from your environment at print time, not baked
into the executable.

\begin{figure}[ht]
\centering
\JarvisVersionBanner
\caption[The \clih{Jarvis --version} banner.]{The \clih{Jarvis --version} banner,
reproduced in the project's own palette; the version on the last line is read from your
installed package at runtime, and here shows the release this manual documents.}
\label{fig:version-banner}
\end{figure}

\begin{jwarn}
If your shell reports \cli{Jarvis: command not found}, the executable is not on your
\code{PATH} --- almost always because you installed into a different environment than
the one you are using now. Re-activate the environment where you ran \cli{pip install},
and confirm with \cli{which Jarvis}. Under \code{pyenv}, \cli{Jarvis} lives in that
version's \code{bin/} directory.
\end{jwarn}

\section{Install Redis}
\label{sec:redis-install}

Pick one route:

\paragraph{macOS (Homebrew).}
\begin{enumerate}
  \item Install and start the service:
\begin{terminal}
\begin{jcmd}
brew install redis
brew services start redis
\end{jcmd}
\end{terminal}
\end{enumerate}

\paragraph{Linux (Debian/Ubuntu, \code{apt}).}
\begin{enumerate}
  \item Install and enable the service:
\begin{terminal}
\begin{jcmd}
sudo apt-get update
sudo apt-get install -y redis-server
sudo systemctl enable --now redis-server
\end{jcmd}
\end{terminal}
\end{enumerate}

\paragraph{Linux (RHEL/Fedora, \code{dnf}).}
\begin{enumerate}
  \item Install and enable the service:
\begin{terminal}
\begin{jcmd}
sudo dnf install -y redis
sudo systemctl enable --now redis
\end{jcmd}
\end{terminal}
\end{enumerate}

\begin{jnote}
On a shared cluster login node you may not have \code{sudo}. Ask the system
administrator to install the \code{redis-server}/\code{redis} package once for
everyone, or fall back to the Docker route below --- it needs no root privilege beyond
running \code{docker} itself.
\end{jnote}

\paragraph{Any platform (Docker).}
\begin{enumerate}
  \item Run the official image on the default port:
\begin{terminal}
\begin{jcmd}
docker run -d --name jarvis-redis -p 6379:6379 redis:7
\end{jcmd}
\end{terminal}
\end{enumerate}

\paragraph{Then verify --- both routes.}
\begin{enumerate}
  \item Ping the service; it answers \code{PONG}:
\begin{terminal}
\begin{jcmd}
redis-cli ping        # -> PONG
\end{jcmd}
\end{terminal}
\end{enumerate}

\begin{jwarn}
Start Redis \emph{before} launching a scan. If the service is down, \cli{Jarvis run}
fails immediately with a connection error --- it never hangs, but it also never starts.
\end{jwarn}

\section{Verify the installation}
\label{sec:verify}

\begin{enumerate}
  \item Confirm the plugin registries respond --- each prints a non-empty list:
\begin{terminal}
\begin{jcmd}
Jarvis portal man     # file formats (JSON, CSV, SLHA, ...)
Jarvis operas list    # registered operators
\end{jcmd}
\end{terminal}
  \item Smoke-test the real pipeline --- sampler, Redis, a Worker, and archiving ---
        with a tiny self-contained card. Save this as \file{smoke.\allowbreak{}yaml}:
\begin{yamlwin}
Scan:
  name: smoke
EnvReqs: {}
Sampling:
  Method: Random
  Bounds:
    point_number: 20
  Variables:
    - name: x
      distribution: {type: Flat, parameters: {min: 0.0, max: 1.0}}
    - name: y
      distribution: {type: Flat, parameters: {min: 0.0, max: 1.0}}

Operas:
  Modules:
    - name: TrivialEggbox
      operator: jarvishep2.testing.eggbox.eggbox2d_numpy
      call_mode: call
      input:
        - {name: x, expression: x}
        - {name: y, expression: y}
      output:
        - {name: z, entry: z}
\end{yamlwin}
        then run it as a smoke test --- \cli{check} draws a handful of points instead
        of the full budget (Chapter~\ref{ch:check-modules}):
\begin{terminal}
\begin{jcmd}
Jarvis check smoke.yaml
\end{jcmd}
\end{terminal}
\end{enumerate}

\yamlkey{Operas.Modules.operator} here is a trivial function shipped inside
\code{jarvishep2} itself --- no external program, no network access, nothing beyond
what you just installed. All three commands succeeded? The installation is complete.

\begin{jnote}
This is deliberately the same card shape Chapter~\ref{ch:quickstart} builds on ---
\cli{check} here just confirms the pipeline works before that chapter walks through it
key by key and runs it for real with \cli{Jarvis run}.
\end{jnote}

\section{Uninstall}
\label{sec:uninstall}

\begin{enumerate}
  \item Remove the package:
\begin{terminal}
\begin{jcmd}
python3 -m pip uninstall Jarvis-HEP
\end{jcmd}
\end{terminal}
\end{enumerate}

Your projects, task cards, and outputs are untouched --- they live in your own
directories, never inside the package (Chapter~\ref{ch:core-concepts}).

\begin{jtip}
Done here? Go straight to Chapter~\ref{ch:quickstart}: it runs a complete 20-point
scan on this installation in under ten minutes.
\end{jtip}
